\begin{thebibliography}{99}
\small\raggedright
\addcontentsline{toc}{section}{References}

\bibitem{AHT2006}
I.~Agol, J.~Hass, and W.~Thurston,
\emph{The computational complexity of knot genus and spanning area},
Transactions of the American Mathematical Society \textbf{358} (2006),
3821--3850.
\url{https://arxiv.org/html/math/0205057v2}.
In particular, Theorems~12 and~16 and Corollary~17.

\bibitem{Tollefson1998}
J.~L.~Tollefson, \emph{Normal surface Q-theory},
Pacific Journal of Mathematics \textbf{183}(2) (1998), 359--374.
\url{https://doi.org/10.2140/pjm.1998.183.359}.

\bibitem{Burton2009}
B.~A.~Burton,
\emph{Converting between quadrilateral and standard solution sets in normal
surface theory}, Algebraic \& Geometric Topology \textbf{9}(4) (2009),
2121--2174.
\url{https://doi.org/10.2140/agt.2009.9.2121};
\url{https://arxiv.org/html/0901.2629v2}.
In particular, Lemma~2.5, Definition~3.9 and Lemma~3.11.

\bibitem{Suh2010}
C.-H.~Suh, \emph{The unknotting problem and normal surface Q-theory},
preprint (2010), arXiv:1009.1500.
\url{https://arxiv.org/abs/1009.1500}.

\bibitem{BurtonOzlen2014}
B.~A.~Burton and M.~Ozlen,
\emph{A fast branching algorithm for unknot recognition with experimental
polynomial-time behaviour}, preprint, arXiv:1211.1079v3 (2014;
originally posted 2012).
\url{https://arxiv.org/html/1211.1079v3}.

\bibitem{LackenbySlides}
M.~Lackenby, \emph{The unknot recognition problem},
Oxford talk slides, March 2021, particularly the displayed
$2^{O((\log n)^3)}$ bound.
\url{https://people.maths.ox.ac.uk/lackenby/quasipolynomial-talk-oxford.pdf}.

\bibitem{Oxford2021}
Mathematical Institute, University of Oxford,
\emph{Marc Lackenby announces a new unknot recognition algorithm that
runs in quasi-polynomial time}, institutional research announcement,
2 February 2021.
\url{https://www.maths.ox.ac.uk/node/38304}.

\bibitem{Lackenby2026}
M.~Lackenby, \emph{Incompressible surfaces, hierarchies and unknot recognition},
preprint, arXiv:2607.23350v1, 25 July 2026.
\url{https://arxiv.org/html/2607.23350v1}.

\bibitem{LackenbyCurves2026}
M.~Lackenby, \emph{Some fast algorithms for curves in surfaces},
Discrete \& Computational Geometry \textbf{76} (2026), 539--588.
\url{https://doi.org/10.1007/s00454-026-00845-7}.
Author text consulted:
\url{https://people.maths.ox.ac.uk/lackenby/AlgorithmsSurfaces13jan26.pdf}.

\bibitem{Regina}
B.~A.~Burton, R.~Budney, W.~Pettersson, et al.,
\emph{Regina: Software for low-dimensional topology}, 1999--2025.
\url{https://regina-normal.github.io/}.
The audit used Python package 7.4.1 (engine reports version 7.4).

\bibitem{Hatcher2002}
A.~Hatcher, \emph{Algebraic topology}, Cambridge University Press, 2002.
Section~3.3, in particular Theorem~3.43 (duality for manifolds with boundary).
\url{https://pi.math.cornell.edu/~hatcher/AT/AT.pdf}.

\bibitem{ProveItSnapshot}
V.~Reshetnikov and ProveIt contributors,
\emph{ProveIt: Topology/UnknotRecognition}, source snapshot
\texttt{7518823550fbc8c217bc8be0113fe52002e77465},
9 October 2026 UTC.
\url{https://github.com/VladimirReshetnikov/ProveIt/tree/7518823550fbc8c217bc8be0113fe52002e77465/Topology/UnknotRecognition}.
The accompanying manifest identifies unchanged baseline dependencies and
the additive continuation files.

\end{thebibliography}
